\chapter{Anatomy}
\label{ch:anatomy}


\section{2026-09-03}

\subsection{Anatomical directional terms}

\begin{itemize}
    \item Superior
    \item Inferior 
    \item Cephalic (towards the head)
    \item Caudal (towards the tail)\footnote{The terms \term{cephalic} and \term{caudal} are often used in embryology, where the head and tail of the embryo are more distinct than in adults.}
    \item Proximal (closer to the point of attachment)
    \item Distal (further from the point of attachment)
    \item Medial
    \item Lateral
    \item Posterior/dorsal
    \item Anterior/ventral
\end{itemize}

\subsection{Anatomical rotational terms}

\begin{itemize}
    \item Supination: rotation of the forearm so that the palm faces anteriorly (or superiorly when the arm is outstretched)
    \item Pronation: rotation of the forearm so that the palm faces posteriorly (or inferiorly when the arm is outstretched)
    \item Inversion: rotation of the foot so that the sole faces medially
    \item Eversion: rotation of the foot so that the sole faces laterally
    \item Flexion: movement that decreases the angle between two body parts\footnote{\label{fn:angle} the `angle' being referred to here is the angle in the `natural' movement (e.g. at the hip it is the raising of the leg that decreases the angle on the anterior side, which is the side that is usually considered when describing flexion and extension).}
    \item Extension: movement that increases the angle between two body parts\footnotemark[\value{footnote}]
    \item Abduction: movement away from the midline of the body
    \item Adduction: movement towards the midline of the body
    \item Circumduction: circular movement of a body part, combining flexion, extension, abduction, and adduction
    \item Rotation: movement around a longitudinal axis, either medially (internal rotation) or laterally (external rotation)
    \item Elevation: movement of a body part superiorly
    \item Depression: movement of a body part inferiorly
    \item Protraction: movement of a body part anteriorly
    \item Retraction: movement of a body part posteriorly
    \item Opposition: movement of the thumb to touch the tips of the other fingers
    \item Reposition: movement of the thumb back to its anatomical position
\end{itemize}
